\documentclass[conference]{IEEEtran}
\IEEEoverridecommandlockouts
\usepackage[UTF8]{ctex}
\usepackage{cite}
\usepackage{amsmath,amssymb}
\usepackage{booktabs}
\usepackage{graphicx}
\usepackage{array}
\usepackage{makecell}
\usepackage{placeins}
\usepackage{url}
\usepackage[hidelinks]{hyperref}

\newcommand{\qwen}{Qwen2.5-Math-1.5B}
\newcommand{\deepseek}{DeepSeek-R1-Distill-Qwen-1.5B}
\newcommand{\mathfive}{MATH-500}

\begin{document}

\title{小型开源大语言模型的数学推理能力：\\提示方法准确率与回答效率评估}

\author{\IEEEauthorblockN{Baizheng Chen}
\IEEEauthorblockA{Department of Computer Science\\
Student ID: 72510800}
}

\maketitle

\begin{abstract}
本文完成 CS6493 课程项目题目一：大语言模型的数学推理能力评估。我们在 \mathfive{}、GSM8K 和 AIME 2024 上评估两个小型数学推理模型：\qwen{} 与 \deepseek{}。方法包括四个基础方法：零样本 Chain-of-Thought、少样本 Chain-of-Thought、Self-Refine 与 Self-Consistency；同时加入两个扩展方法：作为主要纯提示方法的 Plan-and-Solve，以及作为可选工具增强方法的 Tool-Integrated Reasoning。TIR 也可以被解释为轻量 agentic 推理模式，因为模型可以决定何时调用工具、观察结果，并继续完成解题。除准确率和最终回答长度外，本文还报告总生成长度，以及一个基于正确性门控、长度惩罚、答案重复惩罚和过度反思惩罚的行为感知综合分数。
\end{abstract}

\section{引言}
大语言模型能够生成流畅的逐步解答，但数学推理仍然困难，因为模型需要理解符号表达、执行多步逻辑，并输出精确的最终答案。课程项目要求在 \mathfive{}、GSM8K 和 AIME 2024 上比较不同提示方法，并以准确率和回答长度作为关键指标。

本项目在满足要求的基础上加入了效率视角。一个方法如果需要很长、重复的推理才能得到正确答案，并不等价于一个能简洁得到相同答案的方法。这个区别对 Self-Consistency、Self-Refine、Plan-and-Solve 和工具增强方法尤其重要，因为这些方法可能为同一道题触发多轮生成或工具调用。

本文的主要扩展方法是 \texttt{plan\_solve}。它是纯 prompt-only 方法，因此可以与 CoT、Self-Refine 和 Self-Consistency 公平比较。本文也报告 \texttt{tir}，即 Tool-Integrated Reasoning，作为独立的工具增强扩展。由于 TIR 允许 Python 执行，它不被当作普通 prompt-only baseline；但这种设定也是它的主要研究价值，因为它把数学推理推进到 reason、tool invocation、observation、continuation 的 agentic 循环。

\section{相关工作}
CoT prompting 通过显式中间推理步骤提升模型推理能力 \cite{wei2022cot}。Self-Consistency 采样多条推理路径并聚合答案，以提升鲁棒性 \cite{wang2022selfconsistency}。Self-Refine 让模型对自身输出进行批判和修订 \cite{madaan2024selfrefine}。Plan-and-Solve 首先要求模型分解任务，然后按计划求解，主要针对零样本推理中的漏步错误 \cite{wang2023plansolve}。

工具增强方法将自然语言推理与可执行计算结合。ToRA 表明，通过交替生成推理文本和调用外部计算工具，可以显著提升数学求解能力 \cite{gou2023tora}。本文的 TIR 不是训练过的 ToRA 模型，而是一个 prompt 层面的工具循环：当模型认为需要计算时生成 Python 代码块，系统执行后再把结果返回给模型。这使 TIR 不只是更长的提示，而是一个小型 agentic workflow：模型判断是否需要工具，生成可执行 action，接收 observation，再修正剩余推理。

数据集和模型也对应已有研究。GSM8K 测试小学数学文字题 \cite{cobbe2021gsm8k}，MATH 提供竞赛风格数学题 \cite{hendrycks2021math}。Qwen2.5-Math 是面向数学能力优化的 Qwen 模型系列 \cite{yang2024qwenmath}，DeepSeek-R1-Distill-Qwen 是随 DeepSeek-R1 发布的蒸馏推理模型 \cite{deepseek2025r1}。此外，NoWait 研究了 ``Wait'' 等显式反思词带来的推理成本 \cite{wang2025nowait}，Dynamic Early Exit 研究了长推理中的答案感知提前停止 \cite{yang2025earlyexit}。本文只把这些工作作为指标设计动机，而不是声称实现了解码阶段的 early exit。

\section{方法}
\subsection{模型与数据}
本文评估 \qwen{} 和 \deepseek{}。数据集使用项目中的固定采样文件：\mathfive{} 和 GSM8K 各 50 题，AIME 2024 为 30 题。因此每个模型-方法组合共 130 道题。这样的固定抽样设计有利于复现，但也会引入一定的采样偏差。因此，本文报告的各项指标更适合在本项目设定内做相对比较，而不应被理解为对完整数据集真实性能的无偏估计。

\subsection{提示方法}
表~\ref{tab:methods} 定义本文报告的六种方法，并将提示思想和对话形态分开说明。\texttt{cot\_zero} 与 \texttt{cot\_few\_shot} 是单轮方法；\texttt{self\_consistency} 不是迭代式多轮对话，而是对同一题发起五次相互独立的单轮采样；\texttt{self\_refine}、\texttt{plan\_solve} 与 \texttt{tir} 属于多轮或多阶段方法，其中 TIR 还是 agentic 工具循环，因为 Python 执行可以插入模型调用之间，工具输出会成为下一步推理的 observation。

之所以选择这六种方法，是为了覆盖三类不同的干预方式。第一类是 \texttt{cot\_zero} 与 \texttt{cot\_few\_shot}，它们只改变提示上下文，因此是最干净的比较基线。第二类是 \texttt{self\_refine}、\texttt{self\_consistency} 与 \texttt{plan\_solve}，它们通过加入批判、重采样或分解步骤来改变推理内部结构。第三类是 \texttt{tir}，它改变了系统边界本身，因为模型被允许调用外部执行器。这样分类的意义在于：两个方法即使带来相近的 accuracy 提升，背后依赖的额外计算预算也可能完全不同。

把 \texttt{plan\_solve} 选作主要的 prompt-only 扩展方法也是有意为之。与 \texttt{self\_refine} 相比，它是在求解之前增加结构，而不是在求解之后补救，这更符合数学题中“前面漏一步，后面全错”的特点。与 \texttt{self\_consistency} 相比，它不依赖重复采样和投票，因此在有限计算预算下能保持更干净的比较设定。也就是说，Plan-and-Solve 的吸引力不只在于准确率提升，还在于它的机制容易解释，额外代价也更透明。

\texttt{tir} 的角色则不同。它被纳入报告，不是为了替代 prompt-only 比较，而是作为一个对照案例，展示当小模型被允许把算术或符号操作外化给执行器时，会发生什么变化。这一点对本项目很重要，因为数学推理中的失败通常混合了两类来源：一类是解题规划错误，另一类是执行错误。提示分解主要有助于前者，工具调用则可以直接缓解后者。因此同时报告这两种扩展，能够更清楚地区分收益究竟来自更好的 prompting，还是来自改变了整体求解循环。

\begin{table}[t]
\centering
\caption{提示方法定义。}
\label{tab:methods}
\resizebox{\columnwidth}{!}{
\begin{tabular}{llp{0.54\columnwidth}}
\toprule
方法 & 对话形态 & 定义 \\
\midrule
\texttt{cot\_zero} & 单轮 & 零样本 Chain-of-Thought。模型只接收目标题目，被要求逐步推理，并以 boxed final answer 结束。 \\
\texttt{cot\_few\_shot} & 单轮 & 少样本 Chain-of-Thought。提示中加入三个 worked examples，用于约束解题格式和推理风格。 \\
\texttt{self\_refine} & 多轮 & Solve-critique-refine。模型先写初始解答，再批判该解答，最后根据批判重写最终解答。 \\
\texttt{self\_consistency} & 多采样 & 采样投票式 CoT。模型独立采样五个单轮解答，抽取最终答案后多数投票，并选择第一次命中投票答案的完整响应作为最终响应。 \\
\texttt{plan\_solve} & 两轮 & Plan-and-Solve。模型先写简短编号计划，再按该计划完成解题。 \\
\texttt{tir} & Agentic 工具循环 & Tool-Integrated Reasoning。模型可以生成 Python code block，观察返回的 output block，并基于该证据继续推理。 \\
\bottomrule
\end{tabular}}
\end{table}

\subsection{指标}
准确率是主指标。当解析出的预测答案与标准答案在数学意义上相等时，该样本记为正确。回答长度被拆成最终回答长度和总生成长度。最终回答长度是提交给 grader 的文本长度；总生成长度包括中间输出，例如 Plan-and-Solve 的 plan + solve、Self-Refine 的 solve + critique + refine、Self-Consistency 的所有采样，以及 TIR 的完整工具 transcript。当存在完整 token 计数元数据时，总生成 token 使用记录到的生成长度，其中包括生成的代码和记录到的 reasoning 文本；否则 token 长度按 $\lceil\max(\text{字符数}/4,\text{单词数}/0.75)\rceil$ 估算。字符长度是原始字符串长度。

每个样本的综合分数先由正确性门控。效率项定义为：
\begin{equation}
\text{efficiency}_i =
\text{answer}^{0.3}_i
\text{length}^{0.4}_i
\text{reflection}^{0.3}_i.
\end{equation}
最终单样本分数采用 accuracy-first 形式：
\begin{equation}
\text{score}_i = c_i \cdot \left(0.6 + 0.4 \cdot \text{efficiency}_i\right),
\end{equation}
其中 $c_i \in \{0,1\}$ 表示答案是否正确。因此错题仍为 0，答对样本保留主要准确率贡献，效率项只作为次要修正。三个因子的定义如下。

\textbf{长度因子。} 令 $L_i$ 为全流程总生成 token 数，$L_{\text{ref}}=1024$，$\tau_L=512$：
\begin{equation}
\text{length}_i=\exp\left(-\frac{\max(0,L_i-L_{\text{ref}})}{\tau_L}\right).
\end{equation}
该因子只惩罚过长生成，不惩罚短回答。

\textbf{答案因子。} 系统统计答案信号数量 $A_i$，包括 \texttt{\textbackslash boxed\{}、``final answer is''、``the answer is''、``answer:'' 等。尾部比例 $R_i$ 定义为从第一个答案信号到最终响应末尾的字符比例；若没有答案信号，则 $R_i=0$。答案因子为
\begin{align}
\text{count}_i &= \exp\left(-\max(0,A_i-1)\right),\\
\text{tail}_i &= 
\begin{cases}
1, & A_i < 2,\\
\exp\left(-\frac{\max(0,R_i-0.2)}{0.2}\right), & A_i \ge 2,
\end{cases}\\
\text{answer}_i &= \text{count}_i \cdot \text{tail}_i.
\end{align}
它用于惩罚重复声明答案，以及第一次答案之后继续生成过长内容。

\textbf{反思因子。} 系统统计反思提示词数量 $F_i$，例如 ``wait''、``however''、``check''、``verify''、``rethink'' 及中文对应词。参考数量为 1，$\tau_F=2$：
\begin{equation}
\text{reflection}_i=\exp\left(-\frac{\max(0,F_i-1)}{2}\right).
\end{equation}
本文的 joint score 是样本分数均值。正确样本上的平均 efficiency 作为独立诊断指标报告。

这种评分设计有意让 accuracy 保持主导。0.6 的 floor 表示，只要样本答对，大部分分数就已经被保留，效率项只调节剩余的 0.4。这样可以避免一种不合理现象：一个很短但错误的答案，因为“简洁”而被排在一个较长但正确的答案之前。与此同时，次级因子依然有价值，因为数学推理输出即便在 correctness 相同的情况下，也常常在冗长度、自我修正风格和答案呈现方式上差异很大。

这三个因子也分别针对不同的失败模式。长度因子刻画过度生成，尤其对多采样或迭代式方法敏感。答案因子刻画重复给答案、过早给答案后还继续拉长解释等退化行为。反思因子刻画过多犹豫、反复自查，这类现象往往更像推理不稳定，而不是高质量修正。它们都不是要替代 correctness，而是提供一个结构化视角，用来区分那些在 accuracy 上看起来相同、但推理行为其实差异很大的方法。

这种指标设计也回应了一个实际的报告问题。如果只看 accuracy，\texttt{self\_consistency} 会显得特别有吸引力，因为它的投票机制常常能提高正确率，但代价是非常大的隐藏生成成本。如果只看最终回答长度，那么 \texttt{plan\_solve}、\texttt{self\_refine} 和 \texttt{tir} 看起来会比真实情况更“便宜”，因为它们的大部分代价都在中间过程里。如果只看总生成长度，那么稍微冗长但答对的解答又可能被惩罚得过重。因此当前设计是一种折中：首先由 correctness 决定样本是否值得记分，再用效率项区分那些都答对、但代价明显不同的方法。

采用因子分解还有一个重要原因，就是便于误差分析。单个标量只能完成排序，却很难解释为什么两个 accuracy 接近的方法会得到不同的 joint score。因子分解可以。长度因子偏低意味着过度生成，答案因子偏低意味着答案格式不稳定或重复给结论，反思因子偏低则意味着自我修正语言过多。这一点在比较表面上相似的方法时尤其有用，例如 \texttt{cot\_few\_shot} 和 \texttt{plan\_solve}，因为它能帮助判断收益究竟来自更好的任务分解，还是仅仅来自答案风格变化。

\subsection{实验协议}
所有实验都遵循相同的高层流程：先在固定抽样子集上运行推理，再抽取最终答案，计算逐样本指标，最后按模型、方法和数据集聚合结果。所有方法都使用同一套答案抽取器和数学等价性判定器，从而尽量减少仅由格式差异带来的评测漂移。

但不同方法在协议层面仍存在实质差异。\texttt{cot\_zero} 与 \texttt{cot\_few\_shot} 是单次生成，因此它们的 final length 和 total length 相同。\texttt{self\_refine} 会把轨迹扩展成 solve、critique 和 rewrite 三个阶段。\texttt{self\_consistency} 会生成五个独立解答，并对抽取出的答案做多数投票，因此即便最终展示答案很短，它的总生成成本也会显著增加。\texttt{plan\_solve} 显式地把任务分解和执行分开，\texttt{tir} 则把可执行 action 和 observation 引入轨迹。因此，这些方法不能只靠最终 accuracy 比较，还必须同时比较它们为此付出的额外生成代价。

当同一模型、方法和数据集存在多条评测记录时，本文保留 joint score 最高的一条。这种选择是务实的：最终报告希望总结当前项目管线下每种方法最强的一次观测结果。但这也意味着，这些表格更适合被理解为 best-run summary，而不是对多次运行平均方差的估计。

\section{实验}
当同一模型、方法和数据划分存在多条评测记录时，表格只报告 joint score 最高的记录。Overall 行从 overall summary 中独立选择，不由单数据集行重新加权计算。

\begin{table*}[t]
\centering
\caption{130 道题上的 overall 结果。}
\label{tab:overall}
\resizebox{\textwidth}{!}{
\begin{tabular}{llcccc}
\toprule
模型 & 方法 & Accuracy & Final Tok. & Total Tok. & Joint \\
\midrule
\deepseek{} & \texttt{cot\_zero} & 0.785 & 4100.1 & 4100.1 & 0.5112 \\
\deepseek{} & \texttt{cot\_few\_shot} & 0.785 & 3534.8 & 3534.8 & 0.4996 \\
\deepseek{} & \texttt{self\_refine} & 0.792 & 362.3 & 4709.8 & 0.5766 \\
\deepseek{} & \texttt{self\_consistency} & 0.854 & 4217.5 & 21449.8 & 0.5214 \\
\deepseek{} & \texttt{plan\_solve} & 0.892 & 174.1 & 2252.7 & 0.6955 \\
\deepseek{} & \texttt{tir} & 0.908 & 167.5 & 1460.3 & 0.7140 \\
\midrule
\qwen{} & \texttt{cot\_zero} & 0.531 & 2248.2 & 2248.2 & 0.5176 \\
\qwen{} & \texttt{cot\_few\_shot} & 0.515 & 2533.6 & 2533.6 & 0.4604 \\
\qwen{} & \texttt{self\_refine} & 0.508 & 348.5 & 1001.1 & 0.4726 \\
\qwen{} & \texttt{self\_consistency} & 0.562 & 848.9 & 4524.0 & 0.4789 \\
\qwen{} & \texttt{plan\_solve} & 0.785 & 380.9 & 721.7 & 0.6570 \\
\qwen{} & \texttt{tir} & 0.800 & 371.9 & 537.1 & 0.6681 \\
\bottomrule
\end{tabular}}
\end{table*}

\begin{table*}[t]
\centering
\caption{各数据集的 accuracy。}
\label{tab:dataset}
\resizebox{\textwidth}{!}{
\begin{tabular}{llccc}
\toprule
模型 & 方法 & \mathfive{} & GSM8K & AIME 2024 \\
\midrule
\deepseek{} & \texttt{cot\_zero} & 0.900 & 0.900 & 0.400 \\
\deepseek{} & \texttt{cot\_few\_shot} & 0.920 & 0.920 & 0.333 \\
\deepseek{} & \texttt{self\_refine} & 0.880 & 0.920 & 0.433 \\
\deepseek{} & \texttt{self\_consistency} & 0.920 & 0.980 & 0.533 \\
\deepseek{} & \texttt{plan\_solve} & 0.940 & 0.980 & 0.667 \\
\deepseek{} & \texttt{tir} & 0.960 & 0.940 & 0.767 \\
\midrule
\qwen{} & \texttt{cot\_zero} & 0.560 & 0.760 & 0.100 \\
\qwen{} & \texttt{cot\_few\_shot} & 0.560 & 0.780 & 0.000 \\
\qwen{} & \texttt{self\_refine} & 0.560 & 0.740 & 0.033 \\
\qwen{} & \texttt{self\_consistency} & 0.580 & 0.840 & 0.067 \\
\qwen{} & \texttt{plan\_solve} & 0.860 & 0.960 & 0.367 \\
\qwen{} & \texttt{tir} & 0.900 & 0.940 & 0.400 \\
\bottomrule
\end{tabular}}
\end{table*}

\begin{table}[t]
\centering
\caption{正确样本上的 overall 行为因子均值。}
\label{tab:factors}
\resizebox{\columnwidth}{!}{
\begin{tabular}{llccc}
\toprule
模型 & 方法 & Len. & Ans. & Refl. \\
\midrule
DS & \texttt{cot\_zero} & 0.629 & 0.248 & 0.017 \\
DS & \texttt{cot\_few} & 0.727 & 0.083 & 0.039 \\
DS & \texttt{self\_refine} & 0.286 & 0.648 & 0.578 \\
DS & \texttt{self\_cons.} & 0.010 & 0.160 & 0.013 \\
DS & \texttt{plan\_solve} & 0.608 & 0.368 & 0.983 \\
DS & \texttt{tir} & 0.726 & 0.339 & 0.975 \\
\midrule
Qwen & \texttt{cot\_zero} & 1.000 & 0.918 & 0.959 \\
Qwen & \texttt{cot\_few} & 1.000 & 0.585 & 0.961 \\
Qwen & \texttt{self\_refine} & 0.912 & 0.804 & 0.914 \\
Qwen & \texttt{self\_cons.} & 0.515 & 0.758 & 0.942 \\
Qwen & \texttt{plan\_solve} & 0.963 & 0.368 & 0.979 \\
Qwen & \texttt{tir} & 0.972 & 0.368 & 0.955 \\
\bottomrule
\end{tabular}}
\end{table}

表~\ref{tab:overall} 显示，\texttt{plan\_solve} 是两个模型上 accuracy 最强的 prompt-only 扩展方法，并且 joint score 高于四个基线。对 \qwen{}，它把 overall accuracy 从 prompt-only baseline 的 0.531--0.562 提升到 0.785。对 \deepseek{}，更新后的 \texttt{plan\_solve} 达到 0.892 overall accuracy，并在三个数据划分上都保持较强表现。计入完整生成 token 后，它的总生成成本不再像最终回答长度看起来那样低。

TIR 仍然非常有竞争力，在 \deepseek{} 上达到最高 overall accuracy 0.908 和 0.7140 joint score；在 \qwen{} 上达到 0.800 accuracy 和 0.6681 joint score。但 TIR 不能与纯提示方法直接等价比较，因为它允许 Python 执行。因此更合理的解释是：TIR 是一种工具增强且带有 agentic 特征的扩展，展示了当模型可以把部分求解过程外化为可执行 action、观察结果并继续推理时，算术和符号计算如何提升准确率；同时它也会因为代码、工具轮次和记录到的 reasoning 带来额外生成成本。

表~\ref{tab:factors} 说明在 accuracy-first 分数下，效率诊断仍有意义。Self-Consistency 可以有高准确率，但它会采样五个独立输出，因此总生成长度很高。TIR 与 Plan-and-Solve 的最终回答仍较短，但在计入完整生成长度后总 token 成本更高。这说明最终回答长度和总生成成本必须分开报告。

结果还显示出清晰的模型差异。\deepseek{} 在六种方法上都整体强于 \qwen{}，但这种差距并不是均匀的。在两个 proposed 方法下，两个模型都相对于 prompt-only baseline 获得了明显提升，这说明分解式 prompting 和工具调用不是只对某一个 backbone 有效。同时，\deepseek{} 的绝对增益更大，说明更强的原生推理能力更能利用结构化 prompting 与工具交互。

从数据集角度看，GSM8K 对两个模型和大多数方法来说都是最容易的划分，而 AIME 2024 仍然最难。这与直觉一致：基础多步算术更容易直接受益于任务分解和受控执行，而竞赛级题目则需要更深的符号推理。因此，AIME 的结果可以被看作一种 stress test：某些在 GSM8K 上表现良好的方法，未必能顺利迁移到更困难的数学推理场景。

在 prompt-only 方法中，\texttt{plan\_solve} 是最均衡的。它在提升 accuracy 的同时，总生成长度明显低于 \texttt{self\_consistency}，并且通常也低于 \texttt{self\_refine}。这在实践上非常重要。一个方法如果能够增加结构，但不需要显著扩大轨迹数量，就更容易在课程项目中被解释为“仍然是公平的 prompt-based 比较”，同时又带来清晰可见的经验收益。

四个 baseline 也暴露出互补的弱点。\texttt{cot\_zero} 与 \texttt{cot\_few\_shot} 往往生成非常长的回答，尤其是在 \deepseek{} 上，这会在 correctness 尚可时仍拖累效率。\texttt{self\_refine} 虽然能缩短最终回答，但会通过 critique 阶段积累大量隐藏生成成本。\texttt{self\_consistency} 在 baseline 中对 accuracy 的提升最明显，但它的 token 预算也是最高的，因此它恰好说明：如果不把计算成本一起报告，单看 accuracy 很容易得出片面的结论。

\section{讨论}
最稳妥的最终扩展方法是 Plan-and-Solve。它的解释自然：数学解题通常先确定步骤，再执行计算。它也更公平，因为没有使用外部计算器或 verifier。更新后的实验结果显示，它在两个模型上都带来明显 accuracy 收益，但总生成成本需要和简洁的最终回答分开报告。

TIR 更有亮点，但必须谨慎表述。它不应被写成普通 prompt-only baseline，因为它给模型提供了 Python interpreter。它的收益来自语言推理和工具执行的结合。这个设定应该被突出，而不是被弱化：它是一个紧凑的 agentic pattern，模型先选择 action，接收 observation，再更新最终解答。对于数学推理任务，这尤其有价值，因为小模型常见错误正是算术或符号执行错误。它也指向一个比静态 prompting 更宽的方向：小型数学模型可以和可靠工具结合，形成更强的 solver agent。

行为感知指标的意义在于，单纯 accuracy 会隐藏低效推理。Self-Consistency 通过采样提高准确率，但总生成成本高。DeepSeek 的 CoT 运行也出现非常长的最终回答，即使答对也会降低效率。答案因子和反思因子进一步帮助诊断重复答案信号和过度自我修正语言。它们不替代 accuracy，而是作为诊断指标使用。

此外，还需要说明数据层面的限制。由于所有实验都基于固定抽样子集，而不是完整 benchmark，某些增益或下降可能会受样本构成影响。因此，这些表格更适合作为受控设定下的方法比较参考，而不是最终性的绝对结论。

指标设计本身也有第二层限制。本文故意把多种行为压缩进同一个 efficiency 修正项里，这有利于报告简洁，但也会隐藏部分细节。例如，一个方法可能因为稍微过长、重复给答案，或者反思词过多，而得到相近的效率值。因此，每当讨论 joint score 时，因子表都是必要的解释上下文。换句话说，标量分数适合做排序，但真正理解排序原因仍需要依赖因子分解。

最后，一个更方法论层面的结论是：prompt 设计和评测设计应该被一起讨论。方法一旦变成多阶段，最终回答长度就不再是实际推理成本的可靠代理；方法一旦允许调用工具，raw accuracy 也会在缺乏交互设定说明的情况下变得难以解释。因此，本文把方法定义、计算成本核算和评分设计看作一个相互关联的实验选择，而不是三个彼此独立的实现细节。

\section{结论}
本文在 \mathfive{}、GSM8K 与 AIME 2024 上评估了两个小型数学推理模型。最终建议将 Plan-and-Solve 作为主要 proposed prompt-only 方法，因为它公平、简单，且贴合数学解题流程。TIR 可作为可选的工具增强与 agentic 扩展，展示强实验结果，但必须明确说明其 Python 执行设定。最终报告应同时强调 accuracy 与 response efficiency，并区分最终提交答案长度和全流程生成成本。

\FloatBarrier
\bibliographystyle{IEEEtran}
\bibliography{references}

\end{document}
